% SOURCE_AUTHORITY: sga5_fr_workpass.tex, lines 14769--14791 (LF-normalized unit).
% SOURCE_SHA256: 791F4EFFC5E02832D5D77ED03518C8156D6F07E4C8238B03545DB93D883FBB28
Para escrever estas igualdades, identificou-se
\[
\fr_{X\times X'}^*(F\boxt_\Lambda F')
\]
com o produto tensorial $\fr_X^*(F)\otimes_\Lambda\fr_{X'}^*(F')$ mediante o isomorfismo canônico.

Os dois membros da primeira igualdade dependem funtorialmente de $F$ e de $F'$ e são compatíveis com a mudança de base em $X$ ou em $X'$; além disso, são compatíveis com a passagem ao limite indutivo em relação a $F$ ou a $F'$, pois o produto tensorial comuta com o limite indutivo e é claro que
\[
\Fr^*_{(\varinjlim F_i)/X}=\varinjlim\Fr^*_{F_i/X};
\]
Isso permite reduzir ao caso em que $F=\Lambda_X$ e $F'=\Lambda_{X'}$, e então é evidente.

\begin{statement}{Corolário}
Com os dados da proposição 3, os endomorfismos de $\RGamma_{X\times X'}(F\boxt F')$ definidos respectivamente por
\[
(\id_X\times\fr_{X'},\,\id_F\otimes\Fr^*_{F'/X'})
\]
e por
\[
(\fr_X\times\id_{X'},\,\Fr^*_{F/X}\otimes\id_{F'})
\]
são isomorfismos recíprocos entre si.
\end{statement}
